%
% Simulations of air showers and studying observable combinations (chapter 4)
%
\chapter{Combinations of Simulated Air-Shower Observables for Improved Primary Mass Sensitivity}\label{chap:simulations}
%Simulation Dataset Overview. Focus on Auger (simulations with updated hadronic interaction models)

The Heitler-Matthews model, described in~\Cref{chap:airshowers}, presents an overview of how air showers develop in the atmosphere, including the relation of specific air-shower observables with the primary cosmic-ray mass.
Studying these mass-sensitive observables in detail is one method used for mass determination of the primary cosmic ray.
However, after the start of the cascade, it develops in the forward region where particle interactions occur in a phase space with few experimental constraints from terrestrial accelerators and are instead described by collider informed high-energy hadronic interaction models.
Stochasticities in these first few interactions propagate throughout the shower development yielding shower-to-shower fluctuations in air-shower observables for cosmic-ray primaries of the same mass, energy, and arrival direction.
Averaged over many events, certain shower observables are correlated with the cosmic-ray mass, yet fluctuations limit the event-by-event separation for individual observables.
Per-event separation between different primary species can be improved by simultaneously studying multiple mass-sensitive observables.

This chapter presents the results of a mass sensitivity study of individual air-shower observables, and their combination, from cosmic-ray air showers simulated with CORSIKA (COsmic Ray SImulations for KAscade).
All analysis is performed using a general purpose software framework for studying the cosmic-ray mass separation of observables.
The mantra of this framework is to study the intrinsic potential of air-shower observables for event-by-event cosmic-ray mass separation, although CORSIKA simulations need to be computed for location specific conditions such as the elevation above sea level of the observation plane, local geomagnetic field, and atmospheric properties.
With the recently deployed detectors of AugerPrime, the Pierre Auger Observatory simultaneously measures several, complimentary, air-shower observables for primary cosmic-ray mass determination and therefore serves as a promising location choice for the application of this framework.


%Originally, this analysis framework was applied to both the locations of the IceCube Neutrino Observatory and the Pierre Auger Observatory 

%for a set of parameters defined by the characteristics of the CORSIKA simulation library.

%resents a software framework to study the separation power of 

%In preparation for understanding the mass s

%In real analyses, not all detectors will be functioning at the same time, and definitely not all detectors will simultaneously measure high-quality data.

%From the air-shower models described in~\Cref{chap:airshowers},  

%, the reconstructed observables are only statistically related to the mass of the primary cosmic ray.
%On average, there is a correlation between these observables and cosmic-ray mass, yet shower-to-shower fluctuations complicate mass determination at an event level.

%reconstructing cosmic-ray mass, energy, and arrival direction.

%Air-shower universality presented in~\Cref{chap:airshowers} predicts air showers can be fully described using a combination of  ???
\section{CORSIKA Air-Shower Simulations at the Pierre Auger Observatory}\label{sec:simulationdataset}
Within this chapter, a general use CORSIKA simulation library of air showers simulated at the location of the Pierre Auger Observatory is used\footnote[1]{These simulations are available for members of the Pierre Auger Observatory. A brief overview of the library is available at the following link: \href{https://www.auger.unam.mx/AugerWiki/VerticalSimulations78010}{https://www.auger.unam.mx/AugerWiki/VerticalSimulations78010.}}.
The dataset contains CORSIKA v7.8010 air-shower simulations for three separate high-energy hadronic interaction models: EPOS LHC-R, Sibyll 2.3e, and QGSJET III-01.
All low-energy hadronic interactions were modeled with FLUKA v2024.
For each high-energy hadronic interaction model, simulations are available for proton, helium, oxygen, and iron primary cosmic rays over an energy range of $\lg(E / \rm eV)$\,=\,\numrange{14.0}{20.5}, yet only simulations above an energy of $\lg(E / \rm eV)$\,=\,16.0 are studied for this analysis.
Photon primaries are also available in the dataset, but excluded here.
Simulations are separated in energy bins of width $\lg(E / \rm eV)$\,=\,0.5, distributed with an energy spectrum of $dN/dE \propto E^{-1}$ inside the bin resulting in a flat distribution in $\lg(E / \rm eV)$.
Each energy bin contains 5000 air showers per primary\footnote[2]{EPOS LHC-R proton simulations have 10000 air showers per energy bin, but only half were used in this analysis to maintain consistent statistics for each primary}, constituted by 1250 simulations for each of four atmospheric models representing the months of January, March, August, and September in Malarg\"ue, Argentina (in CORSIKA, these respectively are atmospheric models \texttt{MODATM}= 18, 20, 25, 26).

Zenith angles are distributed with $dN/d\theta \propto \sin\theta\cos\theta$ between a range of $\theta$\,=\,\qtyrange{0}{65}{\degree} and azimuth angles are chosen at random.
Table INSERT lists the air-shower simulation statistics for each hadronic model and primary particle used within the context of this analysis.
The total number of simulations available for each hadronic model and primary is approximate as some simulated events contain broken simulation files.
All hadronic models listed in the table will be studied as each model yields differences in the predicted values of the studied mass-sensitive observables that potentially affect their resulting mass sensitivity.  
% put above sentence in following subsection?
% Maybe make this part of the intro???

\subsection{Simulated Air-Shower Observables}\label{sec:observables}
Air-shower observables which can be reconstructed with current ground-based detection techniques were extracted from each CORSIKA simulation.
A parameterized Gaisser-Hillas function, defined by~\Cref{eq:gh_parameterized} in~\Cref{sec:development}, was used to fit the normalized CORSIKA electron-positron longitudinal profile.
Extracted from the fit are the depth of shower maximum ($X_{\rm max}$), the shower width ($L$), and the shower asymmetry ($R$).
Precise air-fluorescence observations of ultra-high-energy cosmic-ray air showers can be used to reconstruct $X_{\rm max}$, $L$, and $R$ (i.e. the full longitudinal profile), while air-shower observations with an ultradense array of radio antennas (such as SKA-low) can reconstruct $X_{\rm max}$ and $L$.
An array of radio antennas with high-quality data, without the need of ultradense spacing, can reconstruct $X_{\rm max}$ alone, with AERA reporting radio $X_{\rm max}$ reconstructions compatible with the Auger fluorescence measurements.
The electron-positron number at shower maximum ($N_{\rm e, max}$) was determined directly from the combined electron-positron CORSIKA longitudinal profile.
$N_{\rm e, max}$ can be reconstructed by either fluorescence telescopes or radio detectors and will be used in this analysis as a proxy to the true air-shower energy for scaling the other observables.

From the CORSIKA particle block, both total muon number at ground ($N_{\mu}$) and the number of electromagnetic particles at ground ($N_{\rm e}$) are obtained.
$N_{\rm e}$ and $N_{\rm e, max}$ correspond only to the combined electron and positron numbers within this study, hence photons are not included in any aspect.
Both $N_{\mu}$ and $N_{\rm e}$ are counted above an energy threshold, where this threshold is intrinsic to the energy limits of tracking separate particle types within the CORSIKA simulations.
These thresholds are \qty{10}{\mega\electronvolt} for muons and \qty{0.25}{\mega\electronvolt} for electrons and positrons within this CORSIKA dataset.
Using the muon and electron numbers at ground, the electron-muon ratio at ground ($\rm R_{\rm e/\mu}$) is calculated.
This ratio has been shown to be a mass-sensitive observable CITEREF5FROMMYPRD, yet requires a particle detector array at ground to have the ability to separate electromagnetic and muonic generated signals.
Discriminating between the two signal types is feasible at Auger from the addition of the SSD.
Refer to~\Cref{chap:airshowers} for theory describing the mass sensitivity of each studied observable. 

In ground-based air shower experiments, the total muon number in an air shower is not directly accessible.
Instead, the muon number is sampled at different distances to the air-shower axis, each serving as a proxy to the total muon content of the air shower. 
In addition to the total muon and electron numbers at ground, these particle numbers are also counted within annuli at different distances to the air-shower axis.
A total of 20 annuli are included, each of \qty{50}{\meter} width, starting from \qtyrange{0}{50}{\meter} from the air-shower axis and ending with an annulus of \qtyrange{950}{1000}{\meter} from the air-shower axis.
Within each annulus, the electron-muon ratio is also calculated.
The annuli are calculated as a distance from the air-shower axis (i.e perpendicular to the direction of the primary cosmic ray) and then projected onto the ground.
This removes the early-late effect in the absolute particle numbers at ground by ensuring particles within a single annulus are sampled from the same shower development region.
The choice of annulus to use for the final mass sensitivity analysis is the \qtyrange{800}{850}{\meter} annulus, referred to as $N_{\mu}$\,\qty{800}{m} and $\rm R_{\rm e/\mu}$\,\qty{800}{\meter} respectively for this annulus' muon number and electron-muon ratio.
This choice of annulus is discussed in~\Cref{app:annulus}. 


\subsubsection{Fitting the Longitudinal Profile of Air Showers}\label{sec:profilefitting}
The particle longitudinal distributions for this CORSIKA simulation library are sampled every \qty{10}{\gram\per\centi\meter\squared} in slant depth.
The combined electron-positron profile is normalized to $N_{\rm e, max}$ and fit as a function of slant depth to~\Cref{eq:gh_parameterized}, where $X^{\prime}$ is replaced with $X - X_{\rm max}$ to enable a direct $X_{\rm max}$ extraction from the fit.
To ensure precise $X_{\rm max}$ values, the normalized particle numbers were weighted with respect to their Poissonian fluctuations, which applied more importance to fitting the longitudinal profile peak.
For all analyses presented within this chapter, only simulated events with fit uncertainties $\sigma_{X_{\rm max}}\,<\,$\qty{5}{\gram\per\centi\meter\squared}, $\sigma_{L}\,<\,$\qty{5}{\gram\per\centi\meter\squared}, and $\sigma_{R}\,<\,$0.05 are included.
This removes from the analysis all simulated air showers with anomalous shower profiles, resulting in the removal of X.X\% of all proton showers and X.X\% of all iron showers.
These are similar to the estimates of anomalous profiles presented in~\cite{Baus:2011kc} and allow for precise knowledge of $X_{\rm max}$, $L$, and $R$ for the simulated events.

%Anomalous shower profiles are typically described by those of double bump showers, air showers where a large portion of the primary energy is transferred to a secondary hadron which travels through the atmosphere before interacting again to initiate a second subshower. 
%Also occurring from this mechanism are elongated air showers with large reconstructed values of $L$.

As shown in XXX, and confirmed in XXX, the $L$ observable can be used in tandem with $X_{\rm max}$ to constrain hadronic interaction models and potentially determine a helium fraction using datasets with large statistics by studying the tail of the $L$ distribution.
To study this further, the full $L$ distributions binned in true energy by $\lg$(E / eV)\,=\,0.5 were investigated.
\Cref{fig:tails} shows the proton and helium $L$ distributions for EPOS LHC-R in a zenith range of \qtyrange{0}{20}{\degree}.
Only proton and helium are included to reduce overcrowding on the plot and because these two mass groups have the largest tails in their respective $L$ distributions.
In this plot, helium consistently has a larger median (thick solid line) and 90\% quantile value (star with thick dotted line) than proton for all energy bins up to \qty[parse-numbers=false]{10^{19.0}}{\electronvolt}, with an up to \qty{10}{\gram\per\centi\meter\squared} difference in the proton-helium 90\% quantile values and differences in the medians of approximately \qty{5}{\gram\per\centi\meter\squared}.
These same results hold for zenith angles $<$\,$\theta_{\rm zen}$\,=\,\qty{40}{\degree} and for the same zenith range using hadronic models QGSJET III-01 and Sibyll 2.3e, yet the exact energy bin where proton transitions to have an equal or larger median and 90\% quantile varies between \qtyrange[parse-numbers=false]{10^{18.0}}{10^{19.0}}{\electronvolt} depending on the hadronic interaction model.
Yet, there is no clear pattern for proton or helium having the longest distribution tail, regardless of energy, zenith, or hadronic model.

\begin{figure}
  \centering
  \includegraphics[trim={0cm 0cm 0cm 0cm},clip,width=1.0\textwidth]{chapter4/Auger_DataCutsApplied_Violin_ProtonHeliumOnly_Lval_Zenith_0_20.pdf}
  \caption[Violin plot of proton-helium $L$ distributions.]{The EPOS LHC-R $L$ distribution for proton and helium primaries in bins of width $\lg$(E / eV) = 0.5. Only zenith angles between \qtyrange{0}{20}{\degree} are included. The median (thick solid line) and 90\% quantile (star with thick dotted line) values are shown for each energy bin.}
  \label{fig:tails}
\end{figure}

Not including the fit quality cuts ($\sigma_{X_{\rm max}}\,<\,$\qty{5}{\gram\per\centi\meter\squared}, $\sigma_{L}\,<\,$\qty{5}{\gram\per\centi\meter\squared}, and $\sigma_{R}\,<\,$0.05) yields similar distributions.
Without cuts, the tails of both the proton and helium distributions become extended, specifically for energies below \qty[parse-numbers=false]{10^{17.0}}{\electronvolt} and above \qty[parse-numbers=false]{10^{19.5}}{\electronvolt}.
Although the difference in the proton-helium 90\% quantiles remains between \qtyrange{5}{10}{\gram\per\centi\meter\squared}, while the medians remain unchanged.
Hence, this further confirms the potential role the $L$ observable tail for discriminating helium primaries from other mass groups if precise reconstructions of $L$ are achieved.

%\subsubsection{Choosing a Muon Annulus for Mass Sensitivity}\label{sec:annulus}

%A single annulus was determined for use in the mass sensitivity analysis by using the figure of merit (FOM) metric, defined as

%\begin{equation}
%    \text{FOM} = \frac{|\mu_{\rm x} - \mu_{\rm y}|}{\sqrt{\sigma_{\rm x}^{2} + \sigma_{\rm y}^{2}}}, \label{eq:fom}
%\end{equation}

%\noindent
%with $\mu_{\rm x/y}$ the mean of distribution x/y and $\sigma_{\rm x/y}$ the x/y distribution's standard deviation.
%A larger FOM represents greater separation between the x and y distributions, assuming they are approximately Gaussian distributed.
%The proton and iron muon 
%to maximize proton-iron separation power


%while attempting to lessen the ratio of $N_{\mu}$ uncertainty-to-counts.
%Due to the distribution of low-energy muons at ground and the Poissonian statistics of surface particle detectors then annuli closer to the shower axis will have smaller relative uncertainties.
%The FOM curves as a function of air-shower energy for the muon number within the farthest 10 annuli from the shower axis are shown in.
%Uncertainties in these FOM curves were estimated from a bootstrapping method, as described in. 

\subsection{Energy Proxy Scaling Corrections for Air-Shower Observables}\label{sec:scaling}

%ensures the particles within an annulus 

%Studying annuli measured in distance from the air-shower axis ensures the particles within an annulus are sampled from the same point of shower development

%hence asymmetries and a zenith dependence on the counted particle numbers are expected.
%As the absolute number and ratios are used, there is no expected asymmetry dependence that propagates through the analysis of the mass sensitivity for these observables; however, there is a zenith dependence.
%This zenith dependence from the air-shower axis calculation is second order in comparison to the zenith angle dependence from total atmospheric depth traversed by the shower.
%This is further discussed in~\Cref{sec:annulus} while describing the choice of annulus used for the mass sensitivity analysis.

%The ground-based observables' zenith dependence from the air-shower axis calculation is second order in comparison to the zenith angle dependence from total atmospheric depth traversed for the shower and therefore is assumed to 

%This ratio has been shown to be a mass-sensitive observable CITEREF5FROMMYPRD, yet requires disentangling signals from electrons in particle detector arrays at ground from those of photons.

%For both the muon and electron numbers at ground, not only are the total particle numbers extracted from CORSIKA, but also are the particle num

%were the total muon number at ground ($N_{\mu}$), the number of electromagnetic particles both at ground and at shower maximum 

%From each CORSIKA simulation, 

%The number of electromagnetic particles at Xmax, Ne;max, can be accessed in experiments, e.g., by fluorescence or radio detectors, and serves as an energy proxy to scale air-shower observables.

%\section{Simulated Air-Shower Observables and their Reconstruction Accuracy}\label{sec:sim_observables}
%\subsection{Fitting the Longitudinal Profile of Air Showers}\label{sec:profilefitting}

\section{Fisher Linear Discriminant Analysis to Model Mass Sensitivity}\label{sec:fisher}
%\subsection{Choosing a Muon Annulus for Mass Sensitivity}\label{sec:annulus}

\section{Gradient Boosted Decision Trees to Improve Mass Sensitivity}\label{sec:gbdt}
Include containment plots